% Kapitel 8 — Budget: 8 Seiten — Schreiben: 14.–23.12.

\chapter{Diskussion}
\label{chap:diskussion}

% Titel am 13.11. angepasst (alt: Antworten auf die Fragen des Exposés und Einordnung), weil die Hauptfrage der Studie (6.1) darin fehlte.
\section[Antworten auf die Hauptfrage der Studie und die Forschungsfragen, Einordnung]{\textcolor{red}{Antworten auf die Hauptfrage der Studie und die Forschungsfragen, Einordnung}}% 4 S.
% Inhalt: FF1 bis FF6 einzeln, auch die Nebenfragen.
\label{sec:antworten}
% NEU 06.11.: Hier die Entwicklung der Fragestellung erzählen (Vorgabe Betreuer 13.10.: FF des Exposés stehen lassen, Wechsel begründen, Entwicklung im Fazit). Reihenfolge: Exposé mit vier Stufen und Hauptgewicht auf Stufe 3 -> AP8 04.09. -> Umstellung auf C gegen D 02.10. -> Stufe 3 gestrichen (Kosten 24.09., Entscheidung nach 19.10.) -> Rechnung 13.10. -> Vorschlag Mail 16.09., Antwort 13.10., Design 19.10. -> Messung 08.10. (Vorher-Nachher). Daten und Quellen in den Notizen von 6.1 (HERKUNFT, Quellenkürzel, 06_nutzerstudie.tex); in der 6.1-Prosa bewusst nicht mehr, dort nur die sachliche Begründung, ausführlich in 5.6.
% NEU 06.11.: Exposé-Kriterium gegen die Annahme der Arbeit (Ausschöpfungsgrad nahe null und kein Vorteil bei den Akzeptanzkriterien; Ex 02_forschungsfragen, letzter Absatz) bei FF2 bis FF4 benennen und sagen, dass es die Arbeit nicht mehr trägt. Aus 6.1 gestrichen, steht sonst nirgends in der Thesis.

\section{Validität und Grenzen}% 4 S.
% Inhalt: Stufe 3 weggelassen, acht Personen, Vortest als Suchauftrag,
% Wertung nicht verblindet, zwei Werkzeugstände.
\label{sec:validitaet}
% =====================================================================
% PFLICHT (13.11.) -- MUSS IN 8.2 STEHEN
% Diese vier Einschränkungen wurden am 13.11. aus Kapitel 6 gestrichen und
% stehen sonst nirgends mehr im Fließtext. Ohne sie fehlen sie in der Arbeit.
% Erst abhaken, wenn jede als ausformulierter Satz in 8.2 steht:
%   [ ] 1. acht Personen (aus 6.1)
%   [ ] 2. Sehen und Wählen nicht trennbar (aus 6.1, Erklärung steht in 6.2)
%   [ ] 3. Übertragbarkeit, ein Team im zweiten Sprint (aus 6.5)
%   [ ] 4. einzelne Agentenläufe, FF4 (aus 6.1)
% Wortlaut wie zuletzt im Text:
% - aus 6.1, acht Personen: "Bei acht Personen lassen die Daten jedoch nicht erkennen, ob ein kleiner Nutzen besteht, kein Nutzen besteht oder der beobachtete Unterschied auf zufällige Schwankungen zurückgeht." In 6.1 steht weiter, dass die Studie keinen Nachweis einer Wirkung führt.
% - aus 6.1, Sehen und Wählen: "Wer kuratiert, hat Code gesehen und zugleich ausgewählt, und die Studie trennt beides nicht." 6.1 verweist nur noch auf 6.2, dort steht die Erklärung.
% - aus 6.5, Übertragbarkeit: "Weil alle acht Teilnehmenden aus einem Team stammen und sich im zweiten Sprint befinden, ist die Übertragbarkeit auf andere Teams und Personen mit anderem Vorwissen zur Codebasis eingeschränkt."
% - aus 6.1, einzelne Agentenläufe (FF4): Aussagekraft der einzelnen Agentenläufe (FF4). Je Person und Aufgabe gibt es einen Lauf mit Prüfliste. Unterschiede zwischen C und D lassen sich damit nicht von Schwankungen zwischen Läufen trennen. Wiederholte automatisierte Läufe mit verschiedenen Arten von Kontext wurden nicht durchgeführt (Zeit, Tokenkosten). Die Arbeit belegt deshalb nicht, ob kuratierter Kontext dem Agenten mehr hilft als automatischer.
% NICHT schreiben: dass wegen der fast vollständigen Deckung des Referenzkontexts beim Ergebnis des Agenten kaum ein Unterschied zu erwarten sei. Die Deckung begründet nur den geringen Spielraum bei der Kontextgüte (6.1, 6.3), nicht das Ergebnis des Agenten.
% Alter Wortlaut aus 6.1 (bis 13.11.): "Die Nutzerstudie untersucht das Ergebnis des Agenten anhand einzelner Läufe und einer Prüfliste. Mehrfach wiederholte Agentenläufe mit verschiedenen Arten von Kontext werden nicht durchgeführt. Sie hätten zusätzliche Zeit und Tokenkosten erfordert. Die bisherigen Messungen zeigten zudem, dass der automatische Abruf auf den untersuchten Anfragen bereits einen großen Teil des Referenzkontexts abdeckt (Abschnitt 6.3). Die Arbeit kann deshalb nicht belegen, ob kuratierter Kontext dem Agenten mehr hilft als automatischer. FF4 wird in der Nutzerstudie anhand der Prüfliste betrachtet, der Vergleich zu FF5 wird nicht untersucht (Tabelle tab:ff-zuordnung)."
% ENDE PFLICHT (13.11.)
% =====================================================================
